%\vspace{\sectionReduceTop}
\section{Introduction}
%\vspace{\sectionReduceBot}

\epigraph{The embodiment hypothesis is the idea that intelligence emerges in the interaction of an agent with
an environment
and as a result of sensorimotor activity.}
{\textit{Smith and Gasser~\cite{smith_al05}}}

\begin{comment}
A key long-term goal of artificial intelligence (AI) is to build intelligent \emph{agents} that can 
%\begin{compactitem}
\begin{inparaitem}[]
    \item \emph{see} (or more generally perceive their environment through vision, audition, or other senses),
    \item \emph{talk} (\ie hold a natural language dialog grounded in the environment),
    %answer natural language questions about its environment),
    \item \emph{act} (\ie navigate their environment and interact with it to accomplish goals), and 
    \item \emph{reason} (\ie consider the long-term consequences of their actions).
%\end{compactitem}
\end{inparaitem}


This goal is of deep scientific and engineering interest. Even a small advance towards such intelligent systems can fundamentally change our lives -- from assistive dialog agents for the visually impaired, to natural-language interaction with self-driving cars, in-home robots, and personal assistants. 
\end{comment}
Imagine walking up to a home robot and asking \myquote{Hey -- can you go check if my laptop is on my desk? And if so, bring it to me.} In order to be successful, such a robot would need a range of skills -- visual perception (to recognize scenes and objects), language understanding (to translate questions and instructions into actions), and navigation in complex environments (to move and find things in a changing environment). 

While there has been significant progress in the vision and language communities thanks to recent advances in deep representations~\cite{he_cvpr16, devlin_arxiv18}, much of this progress has been on `internet AI' rather than \emph{embodied} AI. The focus of the former is pattern recognition in images, videos, and text on \emph{datasets} typically curated from the internet~\cite{Deng09imagenet, mscoco, antol_iccv15}. The focus of the latter is to enable action by an embodied agent (\eg a robot) in an \emph{environment}. This brings to the fore active perception, long-term planning, learning from interaction, and holding a dialog grounded in an environment.

A straightforward proposal is to train agents directly in the physical world -- exposing them to all its richness.
This is valuable and will continue to play an important role in the development of AI.
However, we also recognize that training robots in the real world is 
%\begin{compactitem}
\begin{inparaitem}[]
    \item \emph{slow} (the real world runs no faster than real time and cannot be parallelized), 
    \item \emph{dangerous} (poorly-trained agents can unwittingly injure themselves, the environment, or others), 
    \item \emph{resource intensive} (the robot(s) and the environment(s) in which they execute demand resources and time), 
    \item \emph{difficult to control} (it is hard to test corner-case scenarios as these are, by definition, infrequent and challenging to recreate), and 
    \item \emph{not easily reproducible} (replicating conditions across experiments and institutions is difficult). 
%\end{compactitem}
\end{inparaitem}

%%%%%%%%%%%%%%%%%%%%%%%%%%%%%%%%%%%%%%%%%%%%%%%%%%%%%%%%%%%
\begin{figure*}[t]
    \centering
    %\vspace{-5pt}
    \includegraphics[width=\textwidth]{fig/habitat_stack.pdf}
    %\vspace{-20pt}
    \caption{The `software stack' for training embodied agents involves 
    (1) \emph{datasets} providing 3D assets with semantic annotations, 
    (2) \emph{simulators} that render these assets and 
    within which an embodied agent may be simulated, and 
    (3) \emph{tasks} that define evaluatable problems that enable us to benchmark scientific progress.
    Prior work (highlighted in blue boxes) has contributed a variety of datasets, simulation software, and task definitions.
    We propose a unified embodied agent stack with the Habitat platform, including generic dataset support, a highly performant simulator (\habitatsim), and a flexible API (\habitatapi) allowing the definition and evaluation of a broad set of tasks.}
    %\todo{can we include representative images for all the layers of Habitat stack too?}
    \label{fig:habitat_stack}
    %\vspace{\captionReduceBot}
\end{figure*}
%%%%%%%%%%%%%%%%%%%%%%%%%%%%%%%%%%%%%%%%%%%%%%%%%%%%%%%%%%%

We aim to support a complementary research program: training embodied agents (\eg virtual robots) in rich realistic simulators and then transferring the learned skills to reality.
Simulations have a long and rich history in science and engineering (from aerospace to zoology).
In the context of embodied AI, simulators help overcome the aforementioned challenges -- they can run orders of magnitude faster than real-time and can be parallelized over a cluster; training in simulation is safe, cheap, and enables fair comparison and benchmarking of progress in a concerted community-wide effort.
%
Once a promising approach has been developed and tested in simulation, it can be transferred to physical platforms that operate in the real world~\cite{sim2real-driving,sim2real-legged}.
%In particular, research into mobile robots and other forms of embodied agents can use simulation to evaluate agent policies with standardized benchmarks, establishing a clear measure of progress and promoting reproducibility.

Datasets have been a key driver of progress in computer vision, NLP, and other areas of AI~\cite{Deng09imagenet,mscoco,antol_iccv15,ammirato2017dataset}.
As the community transitions to embodied AI, we believe that simulators will 
assume the role played previously by datasets. 
%Simulators will be used to train and benchmark embodied AI systems that can be reproduced and incrementally advanced by the entire research community. 
To support this transition, 
we aim to standardize the entire `software stack' for training embodied agents 
(\Cref{fig:habitat_stack}): 
%We present \emph{Habitat}, a platform that aims to standardize the entire `software stack' for 
%training embodied agents (\Cref{fig:habitat_stack}): 
scanning the world and creating photorealistic 3D assets, 
developing the next generation of highly efficient and parallelizable simulators, 
specifying embodied AI tasks that enable us to benchmark scientific progress, 
and releasing modular high-level libraries for training and deploying embodied agents. 
%
%\noindent
Specifically, Habitat consists of the following: 
\begin{asparaenum}

\item \habitatsim: 
a flexible, high-performance 3D simulator with configurable agents, 
multiple sensors, and generic 3D dataset handling
(with built-in support for Matterport3D, Gibson, and Replica datasets). 
%and \habitatrep datasets).
%and other datasets). 
%\habitatsim is fast -- 
%when rendering a scene from the Matterport3D dataset, \habitatsim achieves several thousand frames per second (fps) running single-threaded, and can reach over $10@000$ fps multi-process on a single GPU, which is orders of magnitude faster than the closest simulator.

\item \habitatapi: a modular high-level library for 
end-to-end development of embodied AI algorithms -- 
defining embodied AI tasks (\eg navigation,
instruction following, question answering), 
configuring %embodied agents, %(physical form, sensors, capabilities), 
and training embodied agents (via imitation or reinforcement learning, 
or via classic SLAM), and 
benchmarking 
%their performance on the defined tasks 
using standard metrics~\cite{Anderson2018-Evaluation}.

\end{asparaenum}


The Habitat architecture and implementation combine modularity and high performance.
%When rendering a scene from the Matterport3D dataset \habitatsim can reach over $10@000$ fps multi-process on a single GPU, which is orders of magnitude faster than the closest simulator.
When rendering a scene from the Matterport3D dataset, \habitatsim achieves several thousand frames per second (fps) running single-threaded, and can reach over $10@000$ fps multi-process on a single GPU, which is orders of magnitude faster than the closest simulator. 
\habitatapi allows us to train and benchmark embodied agents with different classes of methods and in different 3D scene datasets. 

These large-scale engineering contributions enable us to answer scientific questions requiring experiments that were till now impracticable or `merely' impractical.
Specifically, in the context of point-goal navigation~\cite{Anderson2018-Evaluation}, we 
make two scientific contributions: 
%
\begin{asparaenum}
\item We revisit the comparison between learning and SLAM approaches from two recent works~\cite{mishkin2019benchmarking, kojima2019learn}
and find evidence for the \textbf{opposite conclusion} -- that learning outperforms SLAM if scaled to an order of magnitude more experience than previous investigations.
%if scaled to total experience far surpassing that of previous investigations, and 

\item We conduct the first cross-dataset generalization experiments $\{$train, test$\} \times \{$Matterport3D, Gibson$\}$ for multiple sensors $\{$Blind\footnote{$Blind$ refers to agents with no visual sensory inputs.}, RGB, RGBD, D$\} \times \{$GPS+Compass$\}$ and find that only agents with depth ($D$) sensors generalize well across datasets.
\end{asparaenum}

We hope that our open-source platform and these findings will advance and guide 
future research in embodied AI. 
